% ============================================================================
% INTRODUCTION (REVISED - Democratization-First Framing)
% ============================================================================

\section{Introduction}
\label{sec:introduction}

The transformative potential of Large Language Models (LLMs) remains largely inaccessible to those who need it most. While frontier models like GPT-4o and Claude~3.5~Sonnet demonstrate remarkable capabilities across reasoning, coding, and creative tasks, their deployment costs---\$4--15 per 1,000 queries---create prohibitive barriers for resource-constrained users. Independent researchers cannot afford exploratory experiments. Students developing AI-augmented applications exhaust their budgets after minimal testing. Startups face existential trade-offs between feature richness and runway survival. Even large enterprises struggle to justify LLM integration for high-throughput workflows where inference costs can exceed \$10,000 daily.

This accessibility crisis persists despite a fundamental market evolution: the LLM ecosystem has fragmented into over 80 commercially available models spanning reasoning-optimized generalists (e.g., GPT-4o), domain-specific specialists (e.g., DeepSeek~Coder for programming), and ultra-efficient distillations (e.g., Gemini~Flash~Lite at \$0.10 per 1k queries). \emph{Cheaper alternatives exist}, but practitioners lack the tools to identify when they suffice. The result is systematic over-provisioning: users default to expensive frontier models out of uncertainty, paying 10--100× more than necessary for tasks where specialists deliver equivalent quality.

\paragraph{The Routing Challenge.} What prevents users from simply ``choosing the cheapest model''? The core difficulty is that \emph{model suitability is query-dependent}. A reasoning-heavy mathematical proof requires frontier capabilities, but a structured SQL query generation task is often handled better by compact specialists. Manually routing each query is infeasible at scale. Existing automated routing systems fail to close this accessibility gap for fundamental reasons:

\begin{itemize}[leftmargin=*]
    \item \textbf{Cascading frameworks} (e.g., FrugalGPT~\cite{chen2023frugalgpt}) require manual configuration of model chains and suffer from $O(N)$ latency scaling, making them impractical for the 80+ model ecosystem.
    \item \textbf{Static classifiers} (e.g., RouteLLM~\cite{ong2024routellm}) cannot adapt to evolving model capabilities and require expensive retraining when new models launch weekly.
    \item \textbf{Online learning systems} face prohibitive cold-start costs: standard contextual bandits require 1,000--3,000 queries (\$50--200) to converge, pricing out precisely the users who need cost optimization most urgently.
\end{itemize}

\paragraph{Our Contribution: BanditGPT as an Accessibility Tool.} We present \ours{}, an open-source routing framework designed to democratize access to LLM capabilities through three core principles:

\begin{enumerate}[leftmargin=*]
    \item \textbf{Immediate Usability:} Through \emph{shippable priors}---compact covariance matrices distilled offline from expert supervision---\ours{} achieves 96--99\% reduction in cold-start exploration costs. Users gain production-ready routing from the first query, eliminating the burn-in barrier.
    \item \textbf{Tunable Cost-Quality Trade-Offs:} A single sensitivity parameter ($\lambda$) enables dynamic navigation between ``Cost Leader'' mode (61\% cheaper than FrugalGPT at \$0.70 per 1k queries) and ``High Assurance'' mode (matching frontier-model reliability). This allows students to maximize experimentation volume while enterprises can dial up reliability for production workloads---all without architectural changes.
    \item \textbf{Future-Proof Scalability:} By decoupling prompt-difficulty learning from model-capability profiling, \ours{} supports $O(1)$ onboarding of new models via public metadata. When a cheaper model launches, users benefit immediately without manual recalibration.
\end{enumerate}

\paragraph{Impact Across User Segments.} This framework unlocks previously infeasible use cases:

\begin{itemize}[leftmargin=*]
    \item \textbf{Students \& Educators:} Enable AI-augmented coursework and research projects at \$0.70 per 1k queries (vs.\ \$4.38 for GPT-4o-only), expanding access from elite institutions to under-resourced programs.
    \item \textbf{Independent Researchers:} Afford large-scale qualitative analysis, synthetic data generation, and iterative prompt development that frontier-only pricing makes prohibitive.
    \item \textbf{Startups:} Reduce inference costs by 61--68\%, converting AI features from loss leaders into viable product differentiators while extending runway.
    \item \textbf{Enterprises:} Deploy LLM-powered automation at scale (customer support, document processing) where current costs are blocking adoption despite clear ROI potential.
\end{itemize}

\paragraph{Technical Validation for Trust.} For users to adopt cost-optimized routing, they must trust that quality will not catastrophically degrade. This paper provides rigorous empirical proof:

\begin{itemize}[leftmargin=*]
    \item \textbf{Warm-Start Efficiency (\S\ref{sec:rq1}):} Expert-distilled priors achieve 64.6\% regret reduction compared to cold-start initialization, with <100 calibration examples required.
    \item \textbf{Adaptive Plasticity (\S\ref{sec:rq2}):} \ours{} autonomously corrects for model capability drift and teacher bias, ensuring long-term reliability as the model ecosystem evolves.
    \item \textbf{Pareto Frontier (\S\ref{sec:rq3}):} Establishes state-of-the-art cost-quality trade-offs: Standard mode reduces costs by 61\% vs.\ FrugalGPT (\$0.70 vs.\ \$1.78 per 1k); Hybrid mode matches cascading reliability (95\% accuracy) while operating over an 80+ model pool.
    \item \textbf{Production Viability (\S\ref{sec:engineering}):} Routing overhead is 8.94ms P99 (1.1\% of total latency), ensuring cost savings are effectively free.
\end{itemize}

\paragraph{Open-Source Release for Community Impact.} To maximize accessibility, we release \ours{} as open-source software with pre-trained priors, enabling immediate deployment without data collection or calibration costs. By providing a trustworthy, production-ready tool, we aim to shift the default from ``use the expensive model to be safe'' to ``use adaptive routing to be both effective and sustainable.''

The remainder of this paper presents the technical framework enabling these capabilities (\S\ref{sec:method}), experimental validation across production-scale benchmarks (\S\ref{sec:evaluation}), and analysis of how routing decisions unlock cost arbitrage opportunities (\S\ref{sec:qualitative}).

